\section{Medieval Synthesis and a Doctrinal Boundary}

Aquinas's lectures on John distinguish the spiritual nourishment received by
faith from the chapter's later, explicit treatment of sacramental eating,
without setting them against one another. His sacramental theology then
distinguishes Christ's presence, the sacramental sign, fruitful effect, and the
recipient's disposition. These distinctions explain how Augustine's
``believe, and you have eaten'' can coexist with objective Eucharistic realism:
faith is not a substitute for the sacrament, and merely sacramental reception
does not by itself secure the life promised to one who abides in Christ.

The infant in the manger and Christ in the Eucharist are the same Christ, but
the modes of presence differ. At Bethlehem he is visible in his proper bodily
appearances. In the Eucharist the whole Christ is sacramentally present under
the species of bread and wine. Calling the manger an ``altar'' may be a later
devotional analogy; it is not an identity of sacramental modes.

